\section{卷四 90-106}
\subsection{第90篇(摩西)：   90:1-17}
\subsubsection{主的命定——人都歸於塵土1-6}
主啊，你世世代代作我們的居所。
諸山未曾生出，地與世界你未曾造成，從亙古到永遠，你是神。
你使人歸於塵土，說：你們世人要歸回。
在你看來，千年如已過的昨日，又如夜間的一更。
你叫他們如水沖去；他們如睡一覺。早晨，他們如生長的草，
早晨發芽生長，晚上割下枯乾。
\subsubsection{人生勞苦愁煩空虚短暂，因罪孽活在神的震怒之下 7-11}
我們因你的怒氣而消滅，因你的忿怒而驚惶。
你將我們的罪孽擺在你面前，將我們的隱惡擺在你面光之中。
我們經過的日子都在你震怒之下；我們度盡的年歲好像一聲歎息。
我們一生的年日是七十歲，若是強壯可到八十歲；但其中所矜誇的不過是勞苦愁煩，轉眼成空，我們便如飛而去。
誰曉得你怒氣的權勢？誰按著你該受的敬畏曉得你的忿怒呢？
\subsubsection{四求四願 12-17}
\textcolor{red}{求}你指教我們怎樣數算自己的日子，好叫我們得著智慧的心。
耶和華啊，我們要等到幾時呢？\textcolor{red}{求}你轉回，為你的僕人後悔。
\textcolor{red}{求}你使我們早早飽得你的慈愛，好叫我們一生一世歡呼喜樂。
\textcolor{red}{求}你照著你使我們受苦的日子，和我們遭難的年歲，叫我們喜樂。
\textcolor{red}{願}你的作為向你僕人顯現；\textcolor{red}{願}你的榮耀向他們子孫顯明。
\textcolor{red}{願}主─我們神的榮美歸於我們身上。\textcolor{red}{願}你堅立我們手所做的工；我們手所做的工，\textcolor{red}{願}你堅立。

\subsection{第91篇：   91:1-16}
住在至高者隱密處的，必住在全能者的蔭下。
我要論到耶和華說：他是我的避難所，是我的山寨，是我的神，是我所倚靠的。
他必救你脫離捕鳥人的網羅和毒害的瘟疫。
他必用自己的翎毛遮蔽你；你要投靠在他的翅膀底下；他的誠實是大小的盾牌。
你必不怕黑夜的驚駭，或是白日飛的箭，
也不怕黑夜行的瘟疫，或是午間滅人的毒病。
雖有千人仆倒在你旁邊，萬人仆倒在你右邊，這災卻不得臨近你。
你惟親眼觀看，見惡人遭報。

耶和華是我的避難所；你已將至高者當你的居所，
禍患必不臨到你，災害也不挨近你的帳棚。
因他要為你吩咐他的使者，在你行的一切道路上保護你。
他們要用手托著你，免得你的腳碰在石頭上。
你要踹在獅子和虺蛇的身上，踐踏少壯獅子和大蛇。
神說：因為他專心愛我，我就要搭救他；因為他知道我的名，我要把他安置在高處。
他若求告我，我就應允他；他在急難中，我要與他同在；我要搭救他，使他尊貴。
我要使他足享長壽，將我的救恩顯明給他。

\subsection{第92篇(安息日詩歌)：用琴早晚稱謝神的作為, 惡/義人结局 1-15}
\begin{table}[H]
\centering
\begin{tabular}{p{5.75cm}|p{4.400cm}|p{6.8cm}}\hline\hline
\multicolumn{1}{c|}{用琴早晚稱謝神的作為}&\multicolumn{1}{c|}{惡人结局}&\multicolumn{1}{c}{義人结局}\\\hline
稱謝耶和華！歌頌你至高者的名！
用十弦的樂器和瑟，用琴彈幽雅的聲音，
&\multirow{3}{4.6cm}{畜類人不曉得；愚頑人也不明白。
惡人茂盛如草，一切作孽之人發旺的時候，正是他們要滅亡，直到永遠。
惟你─耶和華是至高，直到永遠。
耶和華啊，你的仇敵都要滅亡；一切作孽的也要離散。}
&\multirow{3}{7.2cm}{你卻高舉了我的角，如野牛的角；我是被新油膏了的。
我眼睛看見仇敵遭報；我耳朵聽見那些起來攻擊我的惡人受罰。
義人要發旺如棕樹，生長如利巴嫩的香柏樹。
他們栽於耶和華的殿中，發旺在我們神的院裡。
他們年老的時候仍要結果子，要滿了汁漿而常發青，
好顯明耶和華是正直的。他是我的磐石，在他毫無不義。}\\\cline{1-1}
早晨傳揚你的慈愛；每夜傳揚你的信實。這本為美事。
&&\\\cline{1-1}
因你耶和華藉著你的作為叫我高興，我要因你手的工作歡呼。
耶和華啊，你的工作何其大!你的心思極其深!&&\\\hline
\end{tabular}
\end{table} 

\subsection{第93篇：耶和華作王世界就堅定,法度的確,殿永稱為聖 1-5}
耶和華作王！他以威嚴為衣穿上；耶和華以能力為衣，以能力束腰，世界就堅定，不得動搖。
你的\textcolor{red}{寶座}從太初立定；你從亙古就有。
耶和華啊，大水揚起，大水發聲，波浪澎湃。
耶和華在高處大有能力，勝過諸水的響聲，洋海的大浪。
耶和華啊，你的\textcolor{red}{法度}最的確；你的\textcolor{red}{殿}永稱為聖，是合宜的。

\subsection{第94篇：伸冤的神用律法教訓的人，用慈愛扶助人却剪除作惡的 1-23}
耶和華啊，你是伸冤的神；伸冤的神啊，求你發出光來！
審判世界的主啊，求你挺身而立，使驕傲人受應得的報應！
耶和華啊，惡人誇勝要到幾時呢？要到幾時呢？
他們絮絮叨叨說傲慢的話；一切作孽的人都自己誇張。
耶和華啊，他們強壓你的百姓，苦害你的產業。
他們殺死寡婦和寄居的，又殺害孤兒。
他們說：耶和華必不看見；雅各的神必不思念。

你們民間的畜類人當思想；你們愚頑人到幾時才有智慧呢？
造耳朵的，難道自己不聽見嗎？造眼睛的，難道自己不看見嗎？
管教列邦的，就是叫人得知識的，難道自己不懲治人嗎？
耶和華知道人的意念是虛妄的。
耶和華啊，你所管教、用律法所教訓的人是有福的！
你使他在遭難的日子得享平安；惟有惡人陷在所挖的坑中。
因為耶和華必不丟棄他的百姓，也不離棄他的產業。
審判要轉向公義；心裡正直的，必都隨從。

誰肯為我起來攻擊作惡的？誰肯為我站起抵擋作孽的？
若不是耶和華幫助我，我就住在寂靜之中了。
我正說我失了腳，耶和華啊，那時你的慈愛扶助我。
我心裡多憂多疑，你安慰我，就使我歡樂。

那藉著律例架弄殘害、在位上行奸惡的，豈能與你相交嗎？
他們大家聚集攻擊義人，將無辜的人定為死罪。
但耶和華向來作了我的高臺；我的神作了我投靠的磐石。
他叫他們的罪孽歸到他們身上。他們正在行惡之中，他要剪除他們；耶和華─我們的神要把他們剪除。

\subsection{第95篇：歌唱感謝神為大又要聽他的話，不可硬著心像祖宗試探神 1-11}
\begin{table}[H]
\centering
\begin{tabular}{p{9.5cm}|p{7.25cm}}\hline\hline
\multicolumn{1}{c|}{歌唱感謝神為大}&\multicolumn{1}{c}{聽他的話，不可硬著心像祖宗試探神}\\\hline
來啊，我們要向耶和華歌唱，向拯救我們的磐石歡呼！
我們要來感謝他，用詩歌向他歡呼！
因耶和華為大神，為大王，超乎萬神之上。
地的深處在他手中；山的高峯也屬他。
海洋屬他，是他造的；旱地也是他手造成的。
來啊，我們要屈身敬拜，在造我們的耶和華面前跪下。
因為他是我們的神；我們是他草場的羊，是他手下的民。
&惟願你們今天聽他的話：
你們不可硬著心，像當日在米利巴，就是在曠野的瑪撒。
那時，你們的祖宗試我探我，並且觀看我的作為。
四十年之久，我厭煩那世代，說：這是心裡迷糊的百姓，竟不曉得我的作為！
所以，我在怒中起誓，說：他們斷不可進入我的安息！\\\hline
\end{tabular}
\end{table} 





\subsection{第96篇：稱頌傳揚神的救恩奇事，敬拜按公義審判世界在主 1-13}
你們要向耶和華唱新歌！全地都要向耶和華歌唱！
要向耶和華歌唱，稱頌他的名！天天傳揚他的救恩！
在列邦中述說他的榮耀！在萬民中述說他的奇事！

因耶和華為大，當受極大的讚美；他在萬神之上，當受敬畏。
外邦的神都屬虛無；惟獨耶和華創造諸天。
有尊榮和威嚴在他面前；有能力與華美在他聖所。

民中的萬族啊，你們要將榮耀、能力歸給耶和華，都歸給耶和華！
要將耶和華的名所當得的榮耀歸給他，拿供物來進入他的院宇。
當以聖潔的（為）妝飾敬拜耶和華；全地要在他面前戰抖！

人在列邦中要說：耶和華作王！世界就堅定，不得動搖；他要按公正審判眾民。
願天歡喜，願地快樂！願海和其中所充滿的澎湃！
願田和其中所有的都歡樂！那時，林中的樹木都要在耶和華面前歡呼。
因為他來了，他來要審判全地。他要按公義審判世界，按他的信實審判萬民。

\subsection{第97篇：神作王显出公義和榮耀，事奉偶像必蒙羞，人靠神必歡喜 1-12}
\begin{table}[H]
\centering
\begin{tabular}{p{7.25cm}|p{10cm}}\hline\hline
\multicolumn{1}{c|}{耶和華作王显出公義和榮耀1-6}&\multicolumn{1}{c}{事奉偶像都蒙羞愧，義人靠神歡喜7-12}\\\hline
耶和華作王！願地快樂！願眾海島歡喜！
密雲和幽暗在他的四圍；公義和公平是他寶座的根基。
有烈火在他前頭行，燒滅他四圍的敵人。
他的閃電光照世界，大地看見便震動。
諸山見耶和華的面，就是全地之主的面，便消化如蠟。
諸天表明他的公義；萬民看見他的榮耀。
&願一切事奉雕刻的偶像、靠虛無之神自誇的，都蒙羞愧。萬神哪，你們都當拜他。
耶和華啊，錫安聽見你的判斷就歡喜；猶大的城邑（原文是女子）也都快樂。
因為你─耶和華至高，超乎全地；你被尊崇，遠超萬神之上。
你們愛耶和華的，都當恨惡罪惡；他保護聖民的性命，搭救他們脫離惡人的手。
散布亮光是為義人；預備喜樂是為正直人。
你們義人當靠耶和華歡喜，稱謝他可記念的聖名。\\\hline
\end{tabular}
\end{table} 




\subsection{第98篇： 两首頌歌-以色列歌唱神的救恩，遍地歌頌神按公義審判 1-9}
\begin{table}[H]
\centering
\begin{tabular}{p{6.7cm}|p{9.6cm}}\hline\hline
\multicolumn{1}{c|}{以色列要歌唱神發明的救恩1-3}&\multicolumn{1}{c}{遍地要歌頌神按公義審判世界4-9}\\\hline
你們要向耶和華唱新歌！因為他行過奇妙的事；他的右手和聖臂施行救恩。
耶和華發明了他的救恩，在列邦人眼前顯出公義；
記念他向以色列家所發的慈愛，所憑的信實。地的四極都看見我們神的救恩。
&全地都要向耶和華歡樂；要發起大聲，歡呼歌頌！
要用琴歌頌耶和華，用琴和詩歌的聲音歌頌他！
用號和角聲，在大君王耶和華面前歡呼！
願海和其中所充滿的澎湃；世界和住在其間的也要發聲。
願大水拍手；願諸山在耶和華面前一同歡呼；
因為他來要審判遍地。他要按公義審判世界，按公正審判萬民。\\\hline
\end{tabular}
\end{table} 

\subsection{第99篇： 神做王坚立公平公义，人要求告并敬拜神 99:1-9}
\begin{table}[H]
\centering
\begin{tabular}{p{5.5cm}|p{5.cm}|p{6.cm}}\hline\hline
\multicolumn{1}{c|}{宣布神做王1-4}&\multicolumn{1}{c|}{神的祭司-求告神5-6}&\multicolumn{1}{c}{神的赎民-遵守神的法度7-9}\\\hline
耶和華作王；萬民當戰抖！他坐在二基路伯上，地當動搖。
耶和華在錫安為大；他超乎萬民之上。
他們當稱讚他大而可畏的名；他本為聖！
王有能力，喜愛公平，堅立公正，在雅各中施行公平和公義。
&你們當尊崇耶和華─我們的神，在他\textcolor{red}{腳凳}前下拜。他\textcolor{red}{本為聖}！
在他的祭司中有摩西和亞倫；在求告他名的人中有撒母耳。他們求告耶和華，他就應允他們。
他在雲柱中對他們說話；
&他們遵守他的法度和他所賜給他們的律例。耶和華─我們的神啊，你應允他們；你是赦免他們的神，卻按他們所行的報應他們。
你們要尊崇耶和華─我們的神，在他的\textcolor{red}{聖山}下拜，因為耶和華─我們的神\textcolor{red}{本為聖}！\\\hline
\end{tabular}
\end{table} 





\subsection{第100篇(称谢诗)：屬他的民当向神歌唱，當讚美進入他的院 1-5}
普天下當向耶和華歡呼！
你們當樂意事奉耶和華，當來向他歌唱！
你們當曉得耶和華是神！我們是他造的，也是屬他的；我們是他的民，也是他草場的羊。
當稱謝進入他的門；當讚美進入他的院。當感謝他，稱頌他的名！
因為耶和華本為善。他的慈愛存到永遠；他的信實直到萬代。


\subsection{第101篇(大卫)： 行完全的道,与完全人同居,遠離惡事，滅絕惡人 1-8}
\subsubsection{歌頌神，立志行完全的道 1-2}
我要歌唱慈愛和公平；耶和華啊，我要向你歌頌！
我要用智慧行完全的道。你幾時到我這裡來呢？我要存完全的心行在我家中。
\subsubsection{遠離惡事，滅絕惡人，与完全人同居 3-8}
邪僻的事，我都不擺在我眼前；悖逆人所做的事，我甚恨惡，不容沾在我身上。
彎曲的心思，我必遠離；一切的惡人（惡事），我不認識。
在暗中讒謗他鄰居的，我必將他滅絕；眼目高傲、心裡驕縱的，我必不容他。
我眼要看國中的誠實人，叫他們與我同住；行為完全的，他要伺候我。
行詭詐的，必不得住在我家裡；說謊話的，必不得立在我眼前。
我每日早晨要滅絕國中所有的惡人，好把一切作孽的從耶和華的城裡剪除。


\subsection{第102篇（大衛\underline{基督的被弃}）： 向神祷告述说苦情，将来要在錫安傳揚讚美神 1-28}
\subsubsection{困苦人發昏的時候向神祷告，并吐露自己的苦情 1-11}
（困苦人發昏的時候，在耶和華面前吐露苦情的禱告。）耶和華啊，求你聽我的禱告，容我的呼求達到你面前！
我在急難的日子，求你向我側耳；不要向我掩面！我呼求的日子，求你快快應允我！
因為，我的年日如煙雲消滅；我的骨頭如火把燒著。
我的心被傷，如草枯乾，甚至我忘記吃飯。
因我唉哼的聲音，我的肉緊貼骨頭。
我如同曠野的鵜鶘；我好像荒場的鴞鳥。
我警醒不睡；我像房頂上孤單的麻雀。
我的仇敵終日辱罵我；向我猖狂的人指著我賭咒。
我吃過爐灰，如同吃飯；我所喝的與眼淚攙雜。
這都因你的惱恨和忿怒；你把我拾起來，又把我摔下去。
我的年日如日影偏斜；我也如草枯乾。
\subsubsection{萬民和列國聚會的時候,在錫安傳揚耶和華的名,并讚美他的話 12-22}
惟你─耶和華必存到永遠；你可記念的名也存到萬代。
你必起來憐恤錫安，因現在是可憐他的時候，日期已經到了。
你的僕人原來喜悅他的石頭，可憐他的塵土。
列國要敬畏耶和華的名；世上諸王都敬畏你的榮耀。
因為耶和華建造了錫安，在他榮耀裡顯現。
他垂聽窮人的禱告，並不藐視他們的祈求。
這必為後代的人記下，將來受造的民要讚美耶和華。
因為，他從至高的聖所垂看；耶和華從天向地觀察，
要垂聽被囚之人的歎息，要釋放將要死的人，
使人在錫安傳揚耶和華的名，在耶路撒冷傳揚讚美他的話，
就是在萬民和列國聚會事奉耶和華的時候。
\subsubsection{你起初立了地的根基,你的年數沒有窮盡 23-28}
他使我的力量中道衰弱，使我的年日短少。
我說：我的神啊，不要使我中年去世。你的年數世世無窮！
你起初立了地的根基；天也是你手所造的。
天地都要滅沒，你卻要長存；天地都要如外衣漸漸舊了。你要將天地如裡衣更換，天地就改變了。
惟有你永不改變；你的年數沒有窮盡。
你僕人的子孫要長存；他們的後裔要堅立在你面前。

\subsection{第103篇（大衛）： 稱頌神的恩惠公義慈愛權柄 1-22}
\subsubsection{我要稱頌神的恩惠 1-5}
\begin{table}[H]
\centering
\begin{tabular}{p{7.25cm}|p{9.75cm}}\hline\hline
\multicolumn{1}{c|}{我的心要稱頌}&\multicolumn{1}{c}{稱頌他的一切恩惠}\\\hline
我的心哪，你要稱頌耶和華！凡在我裡面的，也要稱頌他的聖名！
我的心哪，你要稱頌耶和華！不可忘記他的一切恩惠！
&他赦免你的一切罪孽，醫治你的一切疾病。
他救贖你的命脫離死亡，以仁愛和慈悲為你的冠冕。
他用美物使你所願的得以知足，以致你如鷹返老還童。\\\hline
\end{tabular}
\end{table}

\subsubsection{我們要稱頌神的公義慈愛 6-14}
\begin{table}[H]
\centering
\begin{tabular}{p{7.cm}|p{10.5cm}}\hline\hline
\multicolumn{1}{c|}{耶和華待摩西}&\multicolumn{1}{c}{待我們的慈愛}\\\hline
耶和華施行公義，為一切受屈的人伸冤。
他使摩西知道他的法則，叫以色列人曉得他的作為。
耶和華有憐憫，有恩典，不輕易發怒，且有豐盛的慈愛。
他不長久責備，也不永遠懷怒。
&他沒有按我們的罪過待我們，也沒有照我們的罪孽報應我們。
天離地何等的高，他的慈愛向敬畏他的人也是何等的大！
東離西有多遠，他叫我們的過犯離我們也有多遠！
父親怎樣憐恤他的兒女，耶和華也怎樣憐恤敬畏他的人！
因為他知道我們的本體，思念我們不過是塵土。\\\hline
\end{tabular}
\end{table} 

\subsubsection{你們一切被造的要稱頌神的慈愛權柄 15-22}
\begin{table}[H]
\centering
\begin{tabular}{p{8.cm}|p{3.75cm}|p{4.75cm}}\hline\hline
\multicolumn{1}{c|}{耶和華待世人}&\multicolumn{1}{c}{大能的天使}&\multicolumn{1}{c}{你們作他僕役的}\\\hline
至於世人，他的年日如草一樣。他發旺如野地的花，
經風一吹，便歸無有；他的原處也不再認識他。
但耶和華的慈愛歸於敬畏他的人，從亙古到永遠；他的公義也歸於子子孫孫
就是那些遵守他的約、記念他的訓詞而遵行的人。
&耶和華在天上立定寶座；他的權柄（原文是國）統管萬有。
聽從他命令、成全他旨意、有大能的天使，都要稱頌耶和華！
&你們作他的諸軍，作他的僕役，行他所喜悅的，都要稱頌耶和華！
你們一切被他造的，在他所治理的各處，都要稱頌耶和華！我的心哪，你要稱頌耶和華！\\\hline
\end{tabular}
\end{table} 



\subsection{第104篇： 稱頌神的尊榮威嚴创造和神掌管万有 104:1-35}
\begin{itemize}
\itemsep-.45em
\item 神的尊荣和威严 1-4
\item 神掌管大地 5-8
\item 神掌管大水 9-13
\item 掌管植物 14-18
\item 设立节期 19-23
\item 海中的活物 24-30
\item 要一生都赞美神 31-35
\end{itemize}

\subsubsection{神的尊榮威嚴和创造 1-10}
\begin{table}[H]
\centering
\begin{tabular}{p{8.cm}|p{8.75cm}}\hline\hline
\multicolumn{1}{c|}{神的尊榮威嚴1-4}&\multicolumn{1}{c}{神的创造5-10}\\\hline
我的心哪，你要稱頌耶和華！耶和華─我的神啊，你為至大！你以尊榮威嚴為衣服，
披上亮光，如披外袍，鋪張穹蒼，如鋪幔子，
在水中立樓閣的棟梁，用雲彩為車輦，藉著風的翅膀而行，
以風為使者，以火焰為僕役，
&將地立在根基上，使地永不動搖。
你用深水遮蓋地面，猶如衣裳；諸水高過山嶺。
你的斥責一發，水便奔逃；你的雷聲一發，水便奔流。
諸山升上，諸谷沉下（或譯：隨山上翻，隨谷下流），歸你為他所安定之地。
你定了界限，使水不能過去，不再轉回遮蓋地面。
耶和華使泉源湧在山谷，流在山間，\\\hline
\end{tabular}
\end{table} 


\subsubsection{神供給人、動物物食物和水，掌管節令生死 11-30}
使野地的走獸有水喝，野驢得解其渴。
天上的飛鳥在水旁住宿，在樹枝上啼叫。
他從樓閣中澆灌山嶺；因他作為的功效，地就豐足。
他使草生長，給六畜吃，使菜蔬發長，供給人用，使人從地裡能得食物，
又得酒能悅人心，得油能潤人面，得糧能養人心。
佳美的樹木，就是利巴嫩的香柏樹，是耶和華所栽種的，都滿了汁漿。
雀鳥在其上搭窩；至於鶴，松樹是他的房屋。
高山為野山羊的住所；巖石為沙番的藏處。
你安置月亮為定節令；日頭自知沉落。
你造黑暗為夜，林中的百獸就都爬出來。
少壯獅子吼叫，要抓食，向神尋求食物。
日頭一出，獸便躲避，臥在洞裡。


人出去做工，勞碌直到晚上。
耶和華啊，你所造的何其多！都是你用智慧造成的；遍地滿了你的豐富。
那裡有海，又大又廣；其中有無數的動物，大小活物都有。
那裡有船行走，有你所造的鱷魚游泳在其中。
這都仰望你按時給他食物。
你給他們，他們便拾起來；你張手，他們飽得美食。
你掩面，他們便驚惶；你收回他們的氣，他們就死亡，歸於塵土。
你發出你的靈，他們便受造；你使地面更換為新。

\subsubsection{願罪人從世上消滅，我要一生向耶和華唱詩 31-35}
願耶和華的榮耀存到永遠！願耶和華喜悅自己所造的！
他看地，地便震動；他摸山，山就冒煙。
我要一生向耶和華唱詩！我還活的時候，要向我神歌頌！
願他以我的默念為甘甜！我要因耶和華歡喜！
願罪人從世上消滅！願惡人歸於無有！我的心哪，要稱頌耶和華！你們要讚美耶和華（哈利路亞）！

\subsection{第105篇： 回顾神赐福以色列从立約->得地的过程 105:1-45}
\subsubsection{稱謝傳揚尋求神 1-6}
你們要\textcolor{red}{稱謝}耶和華，\textcolor{red}{求告}他的名，在萬民中\textcolor{red}{傳揚}他的作為！
要向他唱詩\textcolor{red}{歌頌}，\textcolor{red}{談論}他一切奇妙的作為！
要以他的聖名\textcolor{red}{誇耀}！\textcolor{red}{尋求}耶和華的人，心中應當歡喜！
要尋求耶和華與他的能力，時常尋求他的面。
他僕人亞伯拉罕的後裔，他所揀選雅各的子孫哪，你們要\textcolor{red}{記念}他奇妙的作為和他的奇事，並他口中的判語。

\subsubsection{與亞伯拉罕立約，向以撒所起誓,向雅各定為律例 7-15}
他是耶和華─我們的神；全地都有他的判斷。
他記念他的約，直到永遠；他所吩咐的話，直到千代─
就是與亞伯拉罕所立的約，向以撒所起的誓。
他又將這約向雅各定為律例，向以色列定為永遠的約，
說：我必將迦南地賜給你，作你產業的分。
當時，他們人丁有限，數目稀少，並且在那地為寄居的。
他們從這邦遊到那邦，從這國行到那國。
他不容什麼人欺負他們，為他們的緣故責備君王，
說：不可難為我受膏的人，也不可惡待我的先知。

\subsubsection{約瑟掌管埃及,雅各在含地寄居 16-25}
他命饑荒降在那地上，將所倚靠的糧食全行斷絕，
在他們以先打發一個人去─約瑟被賣為奴僕。
人用腳鐐傷他的腳；他被鐵鍊捆拘。
耶和華的話試煉他，直等到他所說的應驗了。
王打發人把他解開，就是治理眾民的，把他釋放，
立他作王家之主，掌管他一切所有的，
使他隨意捆綁他的臣宰，將智慧教導他的長老。
以色列也到了埃及；雅各在含地寄居。
耶和華使他的百姓生養眾多，使他們比敵人強盛，
使敵人的心轉去恨他的百姓，並用詭計待他的僕人。

\subsubsection{神借摩西亞倫領百姓歡樂承受列國的地为業 26-45}
他打發他的僕人摩西和他所揀選的亞倫，
在敵人中間顯他的神蹟，在含地顯他的奇事。
他命黑暗，就有黑暗；沒有違背他話的。
他叫埃及的\textcolor{red}{水變為血}，叫他們的魚死了。
在他們的地上以及王宮的內室，青蛙多多滋生。
他說一聲，\textcolor{red}{蒼蠅}就成群而來，並有\textcolor{red}{虱子}進入他們四境。
他給他們降下\textcolor{red}{冰雹}為雨，在他們的地上降下火焰。
他也擊打他們的葡萄樹和無花果樹，毀壞他們境內的樹木。
他說一聲，就有\textcolor{red}{蝗蟲螞蚱}上來，不計其數，
吃盡了他們地上各樣的菜蔬和田地的出產。
他又擊殺他們國內一切的\textcolor{red}{長子}，就是他們強壯時頭生的。
他領自己的百姓帶銀子金子出來；他支派中沒有一個軟弱的。
他們出來的時候，埃及人便歡喜；原來埃及人懼怕他們。

他鋪張雲彩當遮蓋，夜間使火光照。
他們一求，他就使鵪鶉飛來，並用天上的糧食叫他們飽足。
他打開磐石，水就湧出；在乾旱之處，水流成河。
這都因他記念他的聖言和他的僕人亞伯拉罕。
他帶領百姓歡樂而出，帶領選民歡呼前往。
他將列國的地賜給他們，他們便承受眾民勞碌得來的，
好使他們遵他的律例，守他的律法。你們要讚美耶和華！


\subsection{第106篇： 回顾百姓的悖逆历史，而赞美神的慈爱 106:1-48}
\begin{itemize}
\itemsep-.45em
\item 赞美神 1-5
\item 回顾百姓在去迦南路上的悖逆的罪 6-33
\item 回顾百姓在应许之地犯罪得罪神 34-48
\end{itemize}

\subsubsection{你們要讚美神,求神使我與你的產業一同誇耀 1-5}
你們要讚美耶和華！要稱謝耶和華，因他本為善；他的慈愛永遠長存！
誰能傳說耶和華的大能？誰能表明他一切的美德？
凡遵守公平；常行公義的，這人便為有福！
耶和華啊，你用恩惠待你的百姓；求你也用這恩惠記念我，開你的救恩眷顧我，
使我見你選民的福，樂你國民的樂，與你的產業一同誇耀。

\subsubsection{陳述百姓在紅海行了悖逆 6-12}
我們與我們的祖宗一同犯罪；我們作了孽，行了惡。
我們的祖宗在埃及不明白你的奇事，不記念你豐盛的慈愛，反倒在紅海行了悖逆。
然而，他因自己的名拯救他們，為要彰顯他的大能，
並且斥責紅海，海便乾了；他帶領他們經過深處，如同經過曠野。
他拯救他們脫離恨他們人的手，從仇敵手中救贖他們。
水淹沒他們的敵人，沒有一個存留。
那時，他們才信了他的話，歌唱讚美他。

\subsubsection{陳述百姓在去迦南路上悖逆的罪 13-33}
\begin{table}[H]
\centering
\begin{tabular}{p{2.25cm}|p{2.00cm}|p{4cm}|p{2.5cm}|p{3.00cm}|p{2.00cm}}\hline\hline
\multicolumn{1}{c|}{起慾心}&\multicolumn{1}{c|}{嫉妒摩西}&\multicolumn{1}{c|}{造牛犢}&\multicolumn{1}{c|}{藐視美地}&\multicolumn{1}{c|}{與外邦神連合}&\multicolumn{1}{c}{米利巴水}\\\hline
等不多時，他們就忘了他的作為，不仰望他的指教，
反倒在曠野大起慾心，在荒地試探神。
他將他們所求的賜給他們，卻使他們的心靈軟弱。
&\multirow{1}{2.2cm}{他們又在營中嫉妒摩西和耶和華的聖者亞倫。
地裂開，吞下大坍，掩蓋亞比蘭一黨的人。
有火在他們的黨中發起；有火焰燒毀了惡人。}
&\multirow{1}{4.2cm}{他們在何烈山造了牛犢，叩拜鑄成的像。
如此將他們榮耀的主換為吃草之牛的像，
忘了神─他們的救主；他曾在埃及行大事，
在含地行奇事，在紅海行可畏的事。
所以，他說要滅絕他們；若非有他所揀選的摩西站在當中（破口），使他的忿怒轉消，恐怕他就滅絕他們。}
&\multirow{1}{2.7cm}{他們又藐視那美地，不信他的話，
在自己帳棚內發怨言，不聽耶和華的聲音。
所以，他對他們起誓：必叫他們倒在曠野，
叫他們的後裔倒在列國之中，分散在各地。}
&他們又與巴力‧毘珥連合，且吃了祭死神（人）的物。
他們這樣行，惹耶和華發怒，便有瘟疫流行在他們中間。
那時，非尼哈站起，刑罰惡人，瘟疫這才止息。
那就算為他的義，世世代代，直到永遠。
&他們在米利巴水又叫耶和華發怒，甚至摩西也受了虧損，
是因他們惹動他的靈，摩西（他）用嘴說了急躁的話。
\\\hline
\end{tabular}
\end{table} 



\subsubsection{陳述百姓在迦南地悖逆的罪 34-48}
\begin{table}[H]
\centering
\begin{tabular}{p{2.25cm}|p{6.2cm}|p{3.00cm}|p{5.00cm}}\hline\hline
\multicolumn{1}{c|}{不滅絕外邦人}&\multicolumn{1}{c|}{把自己兒女祭祀鬼魔}&\multicolumn{1}{c|}{屢次設謀背逆}&\multicolumn{1}{c}{被擄掠}\\\hline
他們不照耶和華所吩咐的滅絕外邦人，
反與他們混雜相合，學習他們的行為，
事奉他們的偶像，這就成了自己的網羅，
&把自己的兒女祭祀鬼魔，
流無辜人的血，就是自己兒女的血，把他們祭祀迦南的偶像，那地就被血污穢了。
這樣，他們被自己所做的污穢了，在行為上犯了邪淫。
所以，耶和華的怒氣向他的百姓發作，憎惡他的產業，
將他們交在外邦人的手裡；恨他們的人就轄制他們。
他們的仇敵也欺壓他們，他們就伏在敵人手下。
&他屢次搭救他們，他們卻設謀背逆，因自己的罪孽降為卑下。
然而，他聽見他們哀告的時候，就眷顧他們的急難，
為他們記念他的約，照他豐盛的慈愛後悔。
&他也使他們在凡擄掠他們的人面前蒙憐恤。
耶和華─我們的神啊，求你拯救我們，從外邦中招聚我們，我們好稱讚你的聖名，以讚美你為誇勝。
耶和華─以色列的神是應當稱頌的，從亙古直到永遠。願眾民都說：阿們！你們要讚美耶和華！
\\\hline
\end{tabular}
\end{table} 












